\clearpage

\section{绪论}

\subsection{研究背景及意义}
随着物联网（IoT，Internet of things）技术的快速发展，连接设备的数量呈指数增长，预计未来将有数百亿台设备接入通信网络\cite{zhang2017iot}。这些设备广泛分布于智能家居、工业自动化、智慧城市、环境监测等领域，形成了一个庞大的生态系统\cite{elluri2026exploring}。这一趋势推动了对高效、低功耗、大规模连接通信技术的需求，特别是在设备密集且数据传输量低的场景中尤为突出\cite{2}。针对这种需求，3GPP 提出了海量机器通信（Massive Machine Type Communications，mMTC）作为5G通信的三大应用场景之一\cite{3}，mMTC 专注于支持大量设备的连接，满足小数据包高频次传输的需求\cite{4}。它催生了一系列新的应用，显著改变了各行各业的运作方式。例如，在智慧城市中，mMTC 技术支持了大规模部署的环境传感器网络，实现对空气质量、水资源消耗和垃圾处理的实时监测与优化\cite{5}。在工业自动化领域，mMTC 使得工业设备能够互联互通，促进了生产过程的智能化和效率提升\cite{6}。此外，mMTC 技术在远程医疗中得到了广泛应用，通过联网医疗设备实现患者生理参数的实时监测，为慢性病管理和远程健康指导提供了重要支持\cite{7}。mMTC 作为 5G 通信技术中备受瞩目的典型应用场景，不仅是当今众多科研工作者的研究热点，更是构建实现万物互联的全球物联网的基石\cite{8}。

在 mMTC 场景中，所传输的数据可以大致分为两类：传感类数据和控制类数据\cite{9}。这两类数据由于应用场景的不同，其性能指标关注的重点也各不相同。其中，传感类数据主要用于监测环境或设备状态，例如环境温度、湿度、空气质量、设备运行状态等\cite{daraseliya2025resource}。这类数据的价值取决于其新鲜度，即数据能够反映目标状态的实时性\cite{saha2021minimum}。传统的时延、吞吐量均无法有效刻画信息时效性\cite{kosta2017age}，具体而言，优化吞吐量需要各个传感器采纳较高的采样率，使得网络的通信资源得以被充分利用，即接近满负荷运行。但与此同时，这样的方式会导致网络拥塞，导致数据包在节点队列中经历漫长的等待，使得数据变得老旧过时。即使数据最终能够成功送达，也已沦为丧失应用价值的老旧数据\cite{chen2016age}。而优化时延往往需要保证网络通畅、减少排队与传输等待，因此常采用更稀疏的采样策略，使得发送频率更低，从而减少拥塞，降低排队时长和端到端延迟。尽管这种做法使得零星产生的数据包能够被瞬间传输，但接收端却会因此陷入长期的状态更新盲区，无法持续、连贯地捕获节点的最新动态。换言之就是数据的传输十分迅速，但数据的时效性也无法得到保障\cite{sun2017update}。因此，时延、吞吐量这两个指标均存在局限性，无法有效地刻画信息时效性，因此无法良好的服务于 mMTC 场景。为解决上述问题，信息年龄(Age of Information，AoI)作为一个新型性能评价指标被提出，并得到学界的广泛关注\cite{11}。该指标从信息接收端定义时效性度量标准，其定义为自最近接收到的传感数据生成时刻起至当前时刻的时间间隔。较低的 AoI 意味着数据包更加新鲜，因此接收端能够据此作出更好的决策，从而显著提升工业物联网中状态感知系统的实时响应能力\cite{khodakhah2024balancing}，所以，AoI 是需要重点关注的核心指标之一。

在这种存在大规模设备连接的物联网场景下，由于接入网络的设备数量极为庞大，传统的中央控制方式已经不再适用\cite{9}。因为这种方式需要高效协调每个设备的接入时序和资源分配，然而，设备数的激增导致中央控制的复杂度显著上升，难以满足实时性和可扩展性的要求\cite{berioli2016modern}。正因如此，随机接入方式成为了一种更合理的选择，因为它避免了对全局状态信息的精确掌控，能够在低复杂度的情况下支持大规模设备的动态接入，此外，其分布式特性使得网络具有较好的鲁棒性，能够适应设备加入或退出的动态变化\cite{12}。然而，随机接入方式也存在明显的缺点，即接入冲突。当多个设备在同一时隙尝试接入相同的资源时，会发生冲突，导致数据包传输失败。这种情况下，退避算法被广泛应用，通过调节冲突设备的重传时间来减少冲突和优化资源利用率\cite{kang2010novel}。因此退避算法的设计对于随机接入机制的性能至关重要。现有的退避算法通常以特定的性能指标为优化目标，如最大化网络吞吐量、最小化数据包传输时延或最小化AoI。这些目标虽然适用于某些特定场景，但在实际工业应用中却并不完全符合需求。在高度敏感的工业物联网中，系统性能往往受制于极端恶劣的个别情况，因此它们不仅关注平均指标，还关注设备的违规概率，比如AoI超过预设门限的概率\cite{13}。一个理想的退避算法不仅要关注平均性能的优化，还要保证整个网络中每个设备的公平性，使得所有设备都能及时地更新其数据包。

然而，现有的退避算法在这一点上存在不足。它们往往偏向于优化某些设备的性能，来牺牲其他设备的接入机会。这种不平衡可能导致某些设备的违规概率过高，从而无法满足实际工业场景中的需求。尤其是在任务紧急性分布不均的情况下，这种优化策略可能进一步加剧网络的不公平性，与工业中强调的低违规概率的目标明显不符\cite{wang2023age}。因此，为了更好地满足工业需求，本文提出了一种新的面向 AoI 的随机接入机制，以提高工业系统中的稳定性和可靠性。

\subsection{国内外研究现状}
\subsubsection{面向信息年龄的接入机制研究现状}
物联网的发展带动了诸如环境监测、车联网以及工业自动化等新型业务，在这些场景中，接收端对信息的实时性有着极高的要求，因为它极其依赖周边信息来进行可靠的决策。在无线通信网络中，MAC 协议的设计对发送到接收端的信息新鲜度有着直接的影响，因此，选择一个合理的接入机制是非常重要的。然而，早期的接入机制主要围绕最大化系统吞吐量与最小化传输时延两大指标展开\cite{rom2012multiple}。正如前面所分析的，吞吐量与时延并不能作为一个刻画信息新鲜度的指标，如果在物联网场景下直接套用面向它们所优化的接入机制，会对系统的安全性和稳定性造成巨大的打击。因此，我们有两个问题需要解决：一是找到能刻画信息新鲜度的指标，二是找到合理的接入机制来保证信息的新鲜度。

针对第一个问题，Kaul 等人在2012年针对车载自组织网络的状态更新场景正式提出了AoI这一概念\cite{kaul2012real}。在\cite{kaul2012real}中，AoI 被严格定义为：从接收端的视角出发，自最近接收到的传感数据生成时刻起至当前时刻的时间间隔。相较于时延，AoI 包含了两部分：一是数据包在网络中的逗留时间，二是两个连续到达的有效状态更新之间的时间间隔。AoI 的提出引来了学术界的广泛关注，许多学者也纷纷投身到面向 AoI 的接入机制研究中来。

\begin{enumerate}[wide]
    \item \textbf{从集中式调度到随机接入的演进}
    
    早期围绕 AoI 优化的研究大多聚焦于统筹调度的集中式网络。在集中式调度中，中心节点被假设能够获取网络中所有节点的实时年龄状态、队列信息以及信道状态信息（Channel State Information，CSI），从而在每个时隙做出全局最优的传输决策。最初的集中式调度往往采用直观的贪婪策略，即在每个时隙调度当前 AoI 最大的节点进行传输。然而，研究表明，当不同节点的信道条件或数据包到达率存在异构性时，纯粹的贪婪策略会导致信道较差的节点频繁占用资源却无法成功传输，从而拖累系统整体的年龄性能。为此，Kadota 等人针对无线广播网络提出了基于最大权重策略的集中式调度算法，将节点的信道可靠性与当前 AoI 结合分配权重，以做出最优的传输策略。他们证明了其在最小化平均 AoI 方面的近似最优性，且具有良好的吞吐量表现\cite{kadota2018scheduling}。

    为了进一步在严格的无线资源约束下求得理论最优解，大量学者将 AoI 的集中式调度问题建模为非稳态多臂老虎机（Restless Multi-Armed Bandit, RMAB）问题\cite{sun2017update}。Maatouk 等人利用 Whittle 指数框架，为资源受限的多节点网络推导了低复杂度的指数调度策略，证明了在节点数量趋于无穷时，基于 Whittle 指数的调度方案能够实现 AoI 的极致优化\cite{maatouk2020optimality}。此外，Hsu 等人对具有随机数据包到达特征的集中式调度机制进行了深入的排队论建模，探讨了数据包管理策略与调度算法的联合优化问题\cite{hsu2019scheduling}。Ceran 等人则进一步考虑了混合自动重传请求（HARQ）机制下的 AoI 集中式调度，揭示了资源约束下重传次数与信息新鲜度之间的权衡关系\cite{ceran2019average}。

    尽管上述集中式调度策略在理论上能够达到 AoI 性能的全局最优或近似最优，但其严苛的假设条件在实际网络中难以满足。首先，集中式调度需要基站实时、连续地收集所有节点的当前年龄与队列状态，这在海量 IoT 设备并发接入的大规模网络中会产生极其庞大的信令开销。其次，在许多自组织网络或去中心化的物联网场景中，根本不存在这样一个具备全局视野的中心调度节点\cite{chen2021real}。因此，设计低开销、易扩展、面向 AoI 的分布式随机接入机制，成为了学术界解决大规模 IoT 网络状态更新问题的核心路径与必然演进方向。

    \item \textbf{基于 ALOHA 的随机接入机制研究}
    
    作为最基础、实现最轻量化的分布式接入协议，ALOHA 框架下的 AoI 优化方案在近年来得到了极其广泛的探讨。学者们试图通过修改节点的接入概率或引入不同形式的退避规则，来抑制盲目并发带来的信道碰撞并降低系统年龄。Jiang 等人最早探讨了根据节点当前 AoI 动态分配接入概率的策略，证明了通过优先保障陈旧节点的传输，可以显著降低系统的平均 AoI\cite{jiang2018can}。随后，Zhan 等人进一步针对平均峰值 AoI 的优化，联合设计了节点的基于年龄的接入概率与数据包到达率\cite{zhan2023peak}。此外，部分学者结合博弈论提出了基于年龄的自私接入博弈模型，节点在考虑传输成本的前提下自适应调整接入策略，以期在抑制碰撞与降低 AoI 之间取得平衡\cite{badia2023game}。

    尽管动态概率机制在理论上有效，但 Yavascan 和 Uysal 在文献\cite{yavascan2021analysis}中指出，类似 SAT 的动态机制要求节点始终保持监听模式以估计网络中的活跃节点规模，这带来了极高的能耗，在功耗受限的 IoT 场景中极不现实。为此，他们提出了阈值 ALOHA 机制。该机制规定，节点在成功传输后直接进入休眠，只有当其 AoI 超过某一预设阈值 $T$ 时，才被允许以一个固定的概率参与信道竞争。通过严密的稳态分析，他们证明了通过设定最优的年龄阈值，能够将参与竞争的活跃节点数量大幅压缩，从而显著降低碰撞。

    \item \textbf{基于 CSMA 的随机接入机制研究}
    
    ALOHA 框架本身缺乏物理层面的冲突避免手段，导致其在信道拥塞时仍会引发大量盲目碰撞。为了进一步优化物联网设备的 AoI 性能，学者们将目光转向了具备载波监听能力的 CSMA 机制，通过在发送前进行信道探测，来降低冲突概率。

    Maatouk 等人在\cite{maatouk2020age}中指出，IEEE 802.11 DCF 这种标准的 CSMA 机制的设计初衷是最大化吞吐量并保证节点间的长期公平性。理论推导表明，相比于专门最小化 AoI 的策略，追求最大吞吐量的标准 CSMA 会带来显著的平均 AoI 性能损失。为了解决这一问题，Maatouk 提出了一种睡眠 CSMA 方案。该机制要求节点在每次成功发送数据包后进入一段睡眠期，退出信道竞争，睡眠时长为一个服从指数分布的随机数，通过这种方式将信道资源让渡给其他急需更新的节点。随后，Baiocchi 等人与 Wang 等人进一步将这种基于睡眠的 CSMA 机制拓展到了突发流量、有限缓存以及异构网络场景中，证明了其在实际场景下的有效性\cite{baiocchi2020age}\cite{wang2024analytical}。

    为了进一步在分布式环境中逼近集中式调度的理论性能上限，Tripathi 等人提出了 Fresh CSMA 机制\cite{tripathi2023fresh}。他们认为固定的休眠参数无法适应动态的年龄演化，因此 Fresh CSMA 允许节点根据自身当前的 AoI 随机选择退避计数器。该机制巧妙地在分布式网络中模拟了集中式的最大权重调度策略，使得 AoI 越大的节点获得的退避计数器越小，从而在激烈的信道竞争中占据绝对优势。

    \item \textbf{随机接入机制的研究现状总结}
    
    随着研究的不断发展，面向 AoI 的随机接入机制也在不断完善和进步。然而，纵观研究现状，即使是 ALOHA 框架下最有代表性的阈值 ALOHA，CSMA 框架下最有代表性的 Fresh CSMA 也都具有局限性。阈值 ALOHA 为成功传输的节点增加了静默期，以降低在信道中竞争的活跃节点数量，从而提高了成功发送的概率来降低 AoI。然而，如果节点数量激增，碰撞将不可避免的发生，由于阈值 ALOHA 缺乏碰撞后的处理措施，活跃节点数量将维持在较高的水平，节点难以成功发送，AoI 也无法下降到阈值以下。最终结果是：网络中几乎所有节点的 AoI 都突破了阈值，阈值 ALOHA 彻底失效。而 Fresh CSMA 在节点每次传输数据包后都根据自身的 AoI 选择退避计时器，AoI 越大的节点获得的退避计时器就越小，反之也是如此，这种机制所展现出的性能进一步逼近了集中式调度的理论性能上限。然而，Fresh CSMA 的分析建立在了“忽略信道碰撞”这一理想化的假设之上，假如这种零碰撞环境失效，其性能将遭受巨大的打击，这点在后续的章节中也将得到印证。

    由此可见，现有接入机制的缺陷在于对碰撞这一客观物理现象的妥协与忽视。因此，我们需要一种正视并充分考虑碰撞因素的随机接入机制。除此之外要说明的是，在分析一个接入机制时，仅仅通过仿真是不够的，对其进行数学建模与性能的分析也是重要的环节，因为通过分析得到的理论结果可以更好地为我们提供指导。可以说，接入机制的完善与进步，与性能分析是脱不开干系的，因此研究性能分析的发展过程也是有必要的。
\end{enumerate}

\subsubsection{性能分析的研究现状}
在随机接入机制的设计与评估中，理论性能分析发挥着不可替代的作用。在早期针对 ALOHA 的分析中，学者们多依赖于对全网信道流量的泊松假设，这种宏观概率模型虽然能够推导出吞吐量的理论极限，但无法精细刻画节点在复杂信道环境下的微观演化过程。而在后续普及的 CSMA/CA 接入网络中，节点间的信道状态又形成了极度复杂的强耦合——全网的碰撞概率取决于所有节点的当前退避状态，而个体的状态又受制于全网的碰撞概率。
为了打破这种理论僵局，Bianchi 提出了一种具有里程碑意义的分析框架\cite{bianchi2002performance}。他创新性地假设在稳态条件下，每个节点每次传输遭遇碰撞的概率是一个不受其他节点状态影响的独立常数。基于这一解耦假设，Bianchi 不仅将复杂的协议行为降维成一个可以求解的自洽方程体系，更成为了后续大量MAC层协议进行严谨理论推导的通用基准框架。

但当性能指标的分析重心转移到 AoI 后，系统状态的耦合变得更加棘手。因为 AoI 的演化不仅取决于数据的传输状态，更与连续数据包之间的生成间隔、节点的退避行为以及信道的瞬时竞争环境在时间维度上紧密耦合。从 AoI 的定义来看，在接收端没有新的数据包到达时，AoI 会不断上升，换言之，AoI 的大小与连续两次数据包的成功传输间隔是密不可分的。不妨记这个间隔为 $Y$，实际上，传统求解平均 AoI 的方法里绕不开求解 $Y$ 的期望。按照概率论中的求解思想，我们需要先求解 $Y$ 的概率密度函数（Probability Density Function，PDF），再对其积分以得到期望值。然而，$Y$ 在系统复杂度上升后也将变得难以分析，在多节点的通信网络中，由于碰撞的因素，节点可能要经历若干次传输才能成功发送，而不同协议又会对碰撞采取不同的处理措施，因此 $Y$ 将由随机个随机变量相加得到，这种情况下求解 PDF 是极其困难的。而 PDF 难以求解，自然就难以分析平均 AoI。

为了解决这个问题，Yates 等人率先引入了随机混合系统（Stochastic Hybrid Systems, SHS）来分析多节点网络中的 AoI 演化\cite{yates2018age}。SHS 方法的核心思想是只关注系统的稳定状态，利用 Bianchi 模型的解耦思想，我们可以利用 SHS 对每一个节点的平均 AoI 进行单独的分析。具体而言，它先将网络中各个节点的平均 AoI 记为一个状态向量，接着去寻找会使这个状态发生变化的事件集合，最后将这两部分联立起来，构建一个稳态平衡方程，通过求解这个方程来得到平均 AoI。实际上，这个方程展开来后会变为一组常规的多元一次线性方程组，每一个未知数都对应一个节点的平均 AoI，我们只要求解这个方程组即可得到我们想要的结果，而不再需要考虑求解 $Y$ 的PDF。借助于 SHS 工具，Mollahosseini 等人评估了不同的可变退避机制在时隙 ALOHA 中对平均 AoI 的影响\cite{mollahosseini2024effect}；Maatouk 等人则成功在 CSMA 环境中推导了平均 AoI 的闭式解\cite{maatouk2020age}；而 Tripathi 等人也借此对 Fresh CSMA 的年龄性能进行了严格的理论定界\cite{tripathi2023fresh}。

总体而言，借助 SHS 工具，学术界对平均 AoI 的分析已取得了丰硕的成果。然而，这也是问题所在，因为大多数研究的重心都集中在平均 AoI 的分析上，却鲜有专门针对 AoI 违规概率的分析论文。在工业自动化领域、车联网等新型物联网业务中，系统不仅关心信息有多新鲜，还关注 AoI 维持在一个安全阈值内的概率，即对 AoI 的违规概率有要求\cite{fiems2024age}。实际上，SHS 工具并不适用于求解 AoI 违规概率，在利用它求解平均 AoI 时，我们只关注均值，而不关心数据的离散程度，由此绕开了求解 PDF 的问题，然而在求解违规概率时，我们无法再绕开它了，因为要知道 AoI 大于阈值的概率，需要清楚 AoI 变化的具体分布，自然需要去求解 PDF，然而 PDF 的求解又十分复杂，导致 AoI 违规概率的求解受阻。因此，我们需要一种新的分析方法来解决这个问题，帮助我们分析 AoI 的违规概率。

\subsection{本文主要贡献}
针对现有面向 AoI 的随机接入机制的局限性，以及现有理论分析工具在求解 AoI 违规概率时所面对的技术壁垒，本文提出了一种正视碰撞因素的随机接入机制，并提出一套数学近似分析方法来实现对系统 AoI 违规概率的闭式推导与性能评估。本文的主要工作与贡献可概括为以下四个方面：
\begin{enumerate}[wide]
\item \textbf{提出等待随机接入机制}

针对阈值 ALOHA 机制在高负载下由于缺乏碰撞后的退避措施而导致的“全网违规”失效问题，以及 Fresh CSMA 机制过度依赖“零碰撞”理想假设的脆弱性，本文提出了一种全新的\textbf{等待随机接入机制}。该机制的核心思想是正视信道碰撞的客观存在，在保留现有改进协议“成功传输后主动让步”的基础上，为节点在遭遇碰撞后引入了合理的等待与退避逻辑。通过引入这种确定的碰撞处理措施，本机制可打破高并发场景下的持续碰撞死锁，使网络在非理想的真实碰撞环境中依然能够维持较好的信息更新能力。

\item \textbf{等待 ALOHA 协议的构建与分析}

本文将所提出的等待机制首先应用于基础的等长时隙 ALOHA 框架，构建了等待 ALOHA 协议。在性能分析方面，为了简化复杂的多节点状态，本文引入了解耦假设，利用二维离散时间马尔可夫链对退避计时器和 AoI 的联合变化过程进行了建模，进而推导出了关于系统稳态工作点的非线性自洽方程，并证明了该方程解的存在性与唯一性。接着，本文分析了节点发生连续碰撞次数的几何分布特性，求解出了两次成功传输间隔的一阶与二阶矩，从而得到了平均 AoI 的计算闭式。在求解 AoI 违规概率时，为了绕开复杂的PDF 无穷卷积计算难题，本文采用了一种基于矩匹配的近似策略，利用一个标准分布来拟合更新间隔的概率分布，成功推导出了 AoI 违规概率的计算闭式。

\item \textbf{等待 CSMA 协议的构建与分析}

考虑到 ALOHA 机制缺乏物理层面的碰撞避免手段，本文进一步将等待机制引入到具备载波侦听能力的 CSMA 框架中，设计了等待 CSMA 协议。在进行性能分析时，本文先利用马尔可夫链对系统状态进行建模。在考虑信道冻结特性的基础上，本文推导出了系统稳态工作点的自洽方程，并给出了保证方程解存在且唯一的严格条件。本文充分考虑了 CSMA 中系统时隙长度会随信道状态变化的随机特性，在此基础上计算出了更新间隔的一阶与二阶矩。随后，继续沿用基于矩匹配的近似分析方法，成功推导出了等待 CSMA 机制下的平均 AoI 与 AoI 违规概率的计算闭式。

\item \textbf{机制的仿真验证与理论价值挖掘}

本文通过数值仿真，验证了理论推导闭式的高度准确性，使其可作为可靠的解析工具为后续研究提供助力。同时，通过与多种现有随机接入机制的横向对比评估，本文得出并验证了一个重要的结论：在具有一定规模的通信网络中，适度的主动退避不仅不会加剧信息老化，反而是打破信道拥塞、提升系统整体时效性的有效途径。这一结论为未来海量机器类通信场景下的高可靠 MAC 层协议设计提供了一种务实且稳健的思路。
\end{enumerate}
	
\subsection{论文章节安排}
本文以优化大规模物联网通信场景下的信息时效性为核心目标，针对传统随机接入机制的不足，提出了一种考虑了碰撞因素的随机接入机制，并利用数学工具对其进行了性能分析。全文共分为五个章节，具体的组织结构安排如下：

第一章：绪论。本章首先阐述了本课题的研究背景及意义，指出了在设备海量接入的通信网络中保障信息新鲜度的核心挑战。随后，详细梳理了国内外在面向 AoI 的接入机制、性能分析的研究现状，剖析了现有机制与分析工具所面临的技术局限。最后，总结了本文的主要研究贡献，并给出了论文的整体章节安排。

第二章：相关理论与技术。本章为后续的协议分析奠定了坚实的理论基础。首先，详细介绍了信息年龄及其衍生指标的定义，并推导了一套适用于离散时间系统的通用计算框架。其次，系统地介绍了离散时间马尔可夫链的相关理论与核心性质，为后续章节中对系统稳态接入行为的建模提供了关键的数学工具。

第三章：等待 ALOHA 的提出及分析。本章将所提等待机制引入到时隙 ALOHA 的网络中，设计了等待 ALOHA 协议。我们利用二维离散时间马尔可夫链对节点的演化过程进行了建模，并严格论证了系统稳态自洽方程解的存在性与唯一性。随后，利用几何分布特性推导出更新间隔的一阶与二阶矩，并创新性地引入矩匹配的近似策略，成功绕开了复杂的无穷卷积求解，得出了平均 AoI 与 AoI 违规概率的计算闭式。最后，通过仿真实验验证了理论闭式的准确性，并与现有机制进行了性能对比分析。

第四章：等待 CSMA 的提出及分析。本章进一步将等待机制引入到具备载波侦听能力与变长系统时隙的 CSMA 网络中，设计了等待 CSMA 协议。针对 CSMA 机制的特性，我们利用马尔可夫链对系统进行建模，并推导了考虑微时隙冻结特性的自洽方程。在此基础上，充分考虑变长时隙的随机性并计算出更新间隔的前两阶矩，继续沿用矩匹配近似范式，成功推导出了该机制下时效性指标的计算闭式。本章末尾通过数值仿真与横向对比，揭示了等待 CSMA 协议在应对大规模并发碰撞、保障系统极端可靠性方面的突出优势。

第五章：总结与展望。本章对全文的研究工作进行了系统性总结，归纳了本文所取得的主要成果，客观分析了当前研究中存在的局限性，并对未来该领域的深入研究方向进行了展望。

